% from ega4-21c-fr.tex
\oldpage[IV-4]{328}
complété intégralement clos, et $\mathcal O_{X,y}$ est une
$\mathcal O_{S,s}$-algèbre formellement lisse pour les topologies préadiques
(17.5.3); puisque $Y$ ne contient aucune composante irréductible de $X_s$, on a
$\dim(\mathcal O_{X,y})>\dim(\mathcal O_{S,s})$ (17.5.8). Si $s\neq\eta$, on a
$\dim(\mathcal O_{S,s})\geq1$; si au contraire $s=\eta$, on a par hypothèse
$\dim(\mathcal O_{X,x})\geq2$; donc dans tous les cas
$\dim(\mathcal O_{X,y})\geq2$. En outre, l'idéal premier de $\mathcal O_{X,y}$
correspondant au point $x$ n'est pas contenu dans
$\mathfrak m_s\mathcal O_{X,y}$, puisque $Y$ ne contient aucune composante
irréductible de $X_s$. Enfin, puisque $y$ est fermé dans $X_s$, $k(y)$ est extension
finie de $k(s)$ $(\mathbf I, 6.4.2)$. Dans tous les cas, on peut appliquer à
$A=\mathcal O_{S,s}$, $B=\mathcal O_{X,y}$ et
$B_\mathfrak q=\mathcal O_{X,x}$ le résultat de (21.14.2), ce qui achève la
démonstration.
\end{proof}

\begin{env}[21.14.4]
\label{IV.21.14.4}
\emph{Remarques}. ---

(i) Il se peut que les énoncés (21.14.1) et (21.14.2) restent valables \emph{sans
hypothèse} sur le corps résiduel de $B$. Énoncé avec cette généralité, le résultat
équivaudrait, en vertu de (21.13.15) et (21.13.12.1) au suivant : Soient $A$ un
anneau local noethérien complet, intègre et intégralement clos, $B$ un anneau local
noethérien qui est une $A$-algèbre formellement lisse, tel que
$\dim(B)>\dim(A)$; alors $B$ est parafactoriel.

(ii) Nous verrons au chap.~VI, en employant une technique de « descente finie »
appliquée au morphisme $Y'\to Y=\operatorname{Spec}(A)$, où $Y'$ est le normalisé de
$Y$, que la conclusion de (21.14.1) (ou de (21.14.2)) reste valable en remplaçant
l'hypothèse que $\widehat A$ est intégralement clos par l'hypothèse que
$\widehat A$ est \emph{réduit}, pourvu que $\dim B\geq\dim A+2$. Si on remplace cette
dernière condition par $\dim B\geq\dim A+3$, on peut même supprimer l'hypothèse que
$\widehat A$ est réduit. De même, la conclusion de (21.14.3) reste valable en
remplaçant l'hypothèse que $X$ est normal par l'hypothèse que $X$ est réduit, pourvu
que $\dim(\mathcal O_{X,x})\geq\dim(\mathcal O_{S,f(x)})+2$.

(iii) Au chap.~III, 3\textsuperscript{e} partie, on prouve que si $f:X\to S$ est un
morphisme lisse, $Y$ une partie fermée localement constructible de $X$ telle que,
pour tout $s\in S$, on ait (avec les notations de (21.14.3))
$\operatorname{codim}(Y_s,X_s)\geq3$, alors le couple $(X,Y)$ est parafactoriel
([41], XII, 4.8). Cette conclusion n'est plus valable lorsqu'on suppose seulement que
$\operatorname{codim}(Y_s,X_s)\geq2$ pour tout $s\in S$ et lorsqu'on ne suppose plus
$X$ réduit. Par exemple, soient $k$ un corps, $A=k[T]/(T^2)$, algèbre des nombres
duaux sur $k$, $S=\operatorname{Spec}(A)$,
$X=\operatorname{Spec}(A[T_1,T_2])$ ($T_1$, $T_2$ indéterminées), de sorte que
$X=\mathbb V_S^2$, $Y$ étant la « section nulle » de ce fibré; si $s$ est l'unique
point fermé de $S$, on a $X_s=\mathbb V_k^2$ et $Y_s$ est la « section nulle » de
$X_s$. Pour voir que le couple $(X,Y)$ n'est pas parafactoriel, il suffit de montrer
que l'anneau $B=A[T_1,T_2]_\mathfrak m$, où $\mathfrak m$ est l'idéal
$(T_1)+(T_2)$ qui définit $Y$, n'est pas parafactoriel (21.13.10). Soit
$B_0=B_{\mathrm{red}}$, et notons $U$ et $U_0$ les complémentaires du point fermé
dans $\operatorname{Spec}(B)$ et $\operatorname{Spec}(B_0)$; en raisonnant comme
dans (21.13.9), on a la suite exacte, extension de (21.13.9.1)
\[
  \Gamma(U,\mathcal O_U)^*\longrightarrow\Gamma(U_0,\mathcal O_{U_0})^*
  \longrightarrow\mathrm H^1(U_0,\mathcal O_{U_0})
  \longrightarrow\operatorname{Pic}(U)\longrightarrow\operatorname{Pic}(U_0).
\]
Or, on a $\Gamma(U,\mathcal O_U)=B$, $\Gamma(U_0,\mathcal O_{U_0})=B_0$ puisque
$\operatorname{prof}(B)=\operatorname{prof}(B_0)=2$ (5.10.5). Par ailleurs
$\operatorname{Pic}(U_0)=0$ puisque $B_0$ est factoriel, et
$\mathrm H^1(U_0,\mathcal O_{U_0})\neq0$ [41, 3, Exemple III-1]; comme
l'homomorphisme $B^*\to B_0^*$ est surjectif, on en conclut que
$\operatorname{Pic}(U)\neq0$, donc que $B$ n'est pas parafactoriel.
